\documentclass{urdpl}     % praca w języku polskim

\usepackage[english,polish]{babel}
\usepackage{polski}
\usepackage[utf8]{inputenc}

% dodatkowe pakiety
\usepackage{mathtools}
\usepackage{amsfonts}
\usepackage{amsmath}
\usepackage{amsthm}
\usepackage[hidelinks]{hyperref}
\usepackage{float}
\usepackage{listings}
\usepackage{graphicx}
\usepackage{subcaption}
\usepackage{booktabs}
\usepackage{multirow}
\usepackage{tabularx}
\usepackage{amssymb}
\usepackage{xcolor}
\usepackage{array}
\usepackage{makecell}
\usepackage[flushleft]{threeparttable}
\usepackage{csquotes}
\usepackage{courier}

% ------------------------
% --- < listingi > ---

\lstloadlanguages{TeX}
\renewcommand{\lstlistlistingname}{Spis listingów}
\renewcommand{\lstlistingname}{Listing}

\lstset{
	literate={ą}{{\k{a}}}1
           {ć}{{\'c}}1
           {ę}{{\k{e}}}1
           {ó}{{\'o}}1
           {ń}{{\'n}}1
           {ł}{{\l{}}}1
           {ś}{{\'s}}1
           {ź}{{\'z}}1
           {ż}{{\.z}}1
           {Ą}{{\k{A}}}1
           {Ć}{{\'C}}1
           {Ę}{{\k{E}}}1
           {Ó}{{\'O}}1
           {Ń}{{\'N}}1
           {Ł}{{\L{}}}1
           {Ś}{{\'S}}1
           {Ź}{{\'Z}}1
           {Ż}{{\.Z}}1,
	basicstyle=\footnotesize\ttfamily,
}

% definicja stylu python
\lstdefinestyle{stylePython}{
    language=Python,
    commentstyle=\color{green},
    keywordstyle=\color{blue},
    numberstyle=\tiny\color{gray},
    stringstyle=\color{red},
    basicstyle=\ttfamily\footnotesize,
    breakatwhitespace=false,
    breaklines=true,
    keepspaces=true,
    numbers=left,
    numbersep=5pt,
    showspaces=false,
    showstringspaces=false,
    showtabs=false,
    tabsize=2
}

% definicja stylu bash
\lstdefinestyle{styleBash}{
    language=bash,
    commentstyle=\color{green!50!black}\itshape,
    keywordstyle=\color{blue},
    numberstyle=\tiny\color{gray},
    stringstyle=\color{red},
    basicstyle=\ttfamily\footnotesize,
    breakatwhitespace=false,
    breaklines=true,
    keepspaces=true,
    numbers=left,
    numbersep=5pt,
    showspaces=false,
    showstringspaces=false,
    showtabs=false,
    tabsize=2
}

\definecolor{lightgray}{rgb}{0.9,0.9,0.9}
\definecolor{green}{rgb}{0,0.6,0}
\definecolor{gray}{rgb}{0.5,0.5,0.5}

% ------------------------
\AtBeginDocument{
	\renewcommand{\tablename}{Tabela}
	\renewcommand{\figurename}{Rys.}
    \newcommand{\listingname}{Listing}
}

% --- < tabele > ---
\newcolumntype{C}[1]{>{\hsize=#1\hsize\centering\arraybackslash}X}

%---------------------------------------------------------------------------

\author{Kacper Bułaś, Mateusz Bocak}
\shortauthor{K. Bułaś, M. Bocak}
\noAlbum{127754, 125104}

\titlePL{Klasyfikacja gestów języka migowego ASL przy użyciu Sign Language MNIST}
\titleEN{Laboratory Report: ASL Hand Sign Classification using Sign Language MNIST}

\shorttitlePL{Klasyfikacja ASL przy użyciu Sign Language MNIST}
\shorttitleEN{Lab: ASL Sign Classification with MNIST}

\thesistype{Grupa: 1}

\thesisDone{Prowadzący}
\supervisor{dr Zbigniew Gomółka}

\degreeprogramme{Informatyka}

\date{2025}

\department{Instytut Informatyki}

\faculty{Wydział Nauk Ścisłych i Technicznych}

\setlength{\cftsecnumwidth}{10mm}

%---------------------------------------------------------------------------
\setcounter{secnumdepth}{4}
\brokenpenalty=10000\relax

%---------------------------------------------------------------------------
\begin{document}

\titlepages

\fancypagestyle{plain}
{
    \fancyhf{}
    \renewcommand{\headrulewidth}{0pt}
    \renewcommand{\footrulewidth}{0pt}
}

\setcounter{tocdepth}{2}
\tableofcontents
\clearpage

% dodanie poszczególnych rozdziałów
\include{chapters/chapter001}
\clearpage
\include{chapters/chapter002}
\clearpage
\include{chapters/chapter003}
\clearpage
\chapter{Wnioski i dodatek}
\label{cha:chapter005}

\section{Wnioski}

\begin{enumerate}
    \item \textbf{CNN przewyższa metody klasyczne} na zadaniu klasyfikacji
          znaków ASL: osiąga 92\% dokładności wobec 82\% dla SVC i Lasu
          Losowego. Architektura konwolucyjna jest naturalnie dostosowana
          do danych obrazowych.
    \item \textbf{SVC i Las Losowy osiągają identyczną dokładność} (82\%)
          pomimo różnych strategii uczenia, co sugeruje, że obie metody
          osiągnęły podobną granicę dla tej reprezentacji danych.
    \item \textbf{Jakość inferencji kamerowej zależy od warunków zewnętrznych}
          --- oświetlenia, tła i pozycji dłoni --- bardziej niż od wyboru
          klasyfikatora.
    \item \textbf{MediaPipe odgrywa kluczową rolę} w przygotowaniu wejścia
          modelu; bez etapu detekcji dłoni system nie nadaje się do praktycznego
          zastosowania.
    \item \textbf{Podejście pikselowe ma istotne ograniczenia praktyczne}.
          Reprezentacja oparta na punktach kluczowych mogłaby znacząco
          poprawić odporność systemu na zmienne warunki otoczenia.
    \item \textbf{Modularna architektura projektu} (rozdzielenie warstw danych,
          modeli i inferencji) ułatwia wymianę komponentów i jest dobrą
          praktyką inżynierską w projektach uczenia maszynowego.
\end{enumerate}

\section*{Dodatek: Instrukcja samodzielnej realizacji ćwiczenia}
\addcontentsline{toc}{section}{Dodatek: Instrukcja samodzielnej realizacji ćwiczenia}

Poniższy przewodnik umożliwia samodzielne przeprowadzenie ćwiczenia od
konfiguracji środowiska po uruchomienie inferencji na żywo.

\subsection*{Krok 1: Pobranie repozytorium}

\begin{lstlisting}[style=styleBash, numbers=none, caption={Klonowanie repozytorium}, label=lst:clone]
git clone https://github.com/JustFiesta/mnist-translator.git
cd mnist-translator
\end{lstlisting}

\subsection*{Krok 2: Instalacja środowiska Python}

\begin{lstlisting}[style=styleBash, numbers=none, caption={Konfiguracja środowiska wirtualnego}, label=lst:venv]
# Instalacja Pythona 3.13 przez menedzera uv
uv python install 3.13
uv venv --python 3.13

# Windows PowerShell
.venv\Scripts\Activate.ps1

# Linux / macOS
source .venv/bin/activate

# Instalacja zaleznosci z pliku lock
uv sync
\end{lstlisting}

\subsection*{Krok 3: Pobranie zbioru danych}

\begin{enumerate}
    \item Przejdź na stronę Kaggle:\\
          \url{https://www.kaggle.com/datasets/datamunge/sign-language-mnist/data}
    \item Pobierz pliki \texttt{sign\_mnist\_train.csv} i \texttt{sign\_mnist\_test.csv}.
    \item Utwórz katalog docelowy i umieść w nim pobrane pliki:
\end{enumerate}

\begin{lstlisting}[style=styleBash, numbers=none, caption={Przygotowanie struktury katalogów na dane}, label=lst:dirs]
# Windows PowerShell
New-Item -ItemType Directory -Path data/raw -Force

# Linux / macOS
mkdir -p data/raw

# Skopiuj pobrane pliki CSV do katalogu data/raw/
\end{lstlisting}

\subsection*{Krok 4: Trenowanie modeli}

\begin{lstlisting}[style=styleBash, numbers=none, caption={Trenowanie wszystkich trzech klasyfikatorów}, label=lst:train-all]
# SVC (domyslny, najszybszy)
uv run python main.py train

# Las Losowy
uv run python main.py train --classifier rf --model artifacts/rf_model.pkl

# CNN (trenowanie moze zajac kilka minut)
uv run python main.py train --classifier cnn --model artifacts/cnn_model.keras
\end{lstlisting}

\subsection*{Krok 5: Ewaluacja modeli}

\begin{lstlisting}[style=styleBash, numbers=none, caption={Ewaluacja wytrenowanych modeli}, label=lst:eval-all]
# Ewaluacja SVC
uv run python main.py eval \
    --model artifacts/svc_model.pkl \
    --dataset data/raw/sign_mnist_test.csv

# Ewaluacja Lasu Losowego
uv run python main.py eval \
    --model artifacts/rf_model.pkl \
    --dataset data/raw/sign_mnist_test.csv

# Ewaluacja CNN
uv run python main.py eval \
    --model artifacts/cnn_model.keras \
    --dataset data/raw/sign_mnist_test.csv
\end{lstlisting}

Zanotuj dokładność każdego modelu i porównaj z wynikami z tabeli~\ref{tab:wyniki}.

\subsection*{Krok 6: Inferencja offline (z pliku CSV)}

\begin{lstlisting}[style=styleBash, numbers=none, caption={Inferencja na wybranym wierszu pliku CSV}, label=lst:infer-csv]
uv run python main.py infer \
    --source csv \
    --row 42 \
    --dataset data/raw/sign_mnist_test.csv \
    --model artifacts/svc_model.pkl
\end{lstlisting}

\subsection*{Krok 7: Inferencja na żywo z kamery}

\begin{lstlisting}[style=styleBash, numbers=none, caption={Uruchomienie inferencji w czasie rzeczywistym}, label=lst:infer-live]
uv run python main.py infer \
    --source camera \
    --camera 0 \
    --model artifacts/svc_model.pkl
\end{lstlisting}

\textbf{Uwagi:}
\begin{itemize}
    \item Przy pierwszym uruchomieniu aplikacja automatycznie pobierze plik
          modelu MediaPipe (\texttt{gesture\_recognizer.task}) --- wymagany
          jest dostęp do Internetu.
    \item Naciśnij klawisz \texttt{q}, aby zakończyć działanie okna kamery.
    \item Jeśli kamera nie jest wykrywana pod indeksem 0, spróbuj
          opcji \texttt{-{}-camera 1}.
\end{itemize}

\subsection*{Krok 8: Uruchomienie testów jednostkowych}

\begin{lstlisting}[style=styleBash, numbers=none, caption={Uruchamianie testów pytest}, label=lst:tests]
# Wszystkie testy
uv run pytest

# Testy z raportem pokrycia kodu
uv run pytest --cov=src --cov-report=term-missing
\end{lstlisting}

\subsection*{Pytania kontrolne}

\begin{enumerate}
    \item Dlaczego litery J i Z nie są uwzględnione w zbiorze Sign Language MNIST?
    \item Jaką rolę pełni MediaPipe w potoku inferencji kamerowej?
    \item Co oznacza stratyfikowany podział danych i dlaczego jest ważny
          przy niezbalansowanych klasach?
    \item Jaka jest różnica między parametrami $C$ i $\gamma$ w SVC z jądrem RBF?
    \item Dlaczego CNN osiąga wyższą dokładność niż SVC na danych obrazowych?
    \item W jaki sposób można poprawić odporność systemu na zmienne oświetlenie?
\end{enumerate}

\clearpage
\chapter{Implementacja systemu}
\label{cha:pierwszyDokument}

\section{Architektura logiczna}
\label{sec:architektura_logiczna}

Projekt został podzielony na trzy warstwy:
\begin{itemize}
    \item \texttt{src/data} -- przygotowanie danych i podziały zbioru,
    \item \texttt{src/model} -- uczenie i ewaluacja klasyfikatorów,
    \item \texttt{src/inference} -- predykcja offline i online (kamera).
\end{itemize}

Warstwa wejściowa CLI znajduje się w \texttt{main.py} i deleguje logikę do
modułów w \texttt{src/}. Takie podejście upraszcza testowanie i ogranicza
sprzężenie między interfejsem użytkownika a logiką przetwarzania.

\section{Przetwarzanie danych}
\label{sec:przetwarzanie_danych}

Klasa \texttt{CsvDataset} ładuje plik CSV i wymusza spójność schematu
(obecność kolumny \texttt{label} oraz dokładnie 784 cechy obrazu). Dla każdej
próbki zwracana jest para:
\begin{center}
    \texttt{(feature\_vector: np.ndarray, label: str)}
\end{center}

Wektor cech jest typu \texttt{float32} i przyjmuje wartości w $[0,1]$.
Etykiety numeryczne przekształcane są do liter poprzez mapowanie:
\begin{equation}
    y_{letter} = \mathrm{chr}(\mathrm{ord}('A') + y_{int})
\end{equation}

Budowa macierzy cech realizowana jest przez \texttt{build\_feature\_matrix},
a podział danych przez \texttt{stratified\_split}. Domyślnie stosowany jest
podział 80/10/10 z zachowaniem proporcji klas.

\section{Uczenie modeli}
\label{sec:uczenie_modeli}

Moduł \texttt{src/model/train.py} obsługuje trzy warianty klasyfikatorów:
\begin{itemize}
    \item \textbf{SVC (RBF)} -- domyślny klasyfikator oparty o jądro radialne,
    \item \textbf{Random Forest} -- las losowy o 200 drzewach,
    \item \textbf{CNN (Keras)} -- sieć konwolucyjna dla wejścia 28$\times$28.
\end{itemize}

Modele klasyczne serializowane są przez \texttt{joblib} do formatu
\texttt{.pkl}. Model konwolucyjny zapisywany jest jako \texttt{.keras}
oraz osobny plik mapowania etykiet \texttt{.labels.json}.

\section{Ewaluacja jakości}
\label{sec:ewaluacja_jakosci}

Funkcja \texttt{evaluate} oblicza dokładność globalną oraz raport klasyfikacji
\texttt{classification\_report} biblioteki scikit-learn \cite{sklearn_docs}.
Wyniki zwracane są również jako słownik, co upraszcza dalsze automatyczne
raportowanie.

\section{Inferencja: CSV i kamera}
\label{sec:inferencja_csv_kamera}

Inferencja offline z jednego wiersza CSV realizowana jest przez
\texttt{CsvRowSource}. Inferencja online wykorzystuje klasę \texttt{CameraSource},
która:
\begin{enumerate}
    \item pobiera klatkę z kamery (OpenCV \cite{opencv_docs}),
    \item wykrywa dłoń przez MediaPipe GestureRecognizer \cite{mediapipe_docs},
    \item wyznacza obszar dłoni i dokonuje kadrowania,
    \item przeskalowuje obraz do 28$\times$28 oraz normalizuje piksele,
    \item przekazuje wektor 784 cech do modelu.
\end{enumerate}

W przypadku braku detekcji dłoni moduł zwraca wektor zerowy, dzięki czemu
interfejs inferencji pozostaje stabilny.

\section{Interfejs CLI}
\label{sec:interfejs_cli}

Plik \texttt{main.py} udostępnia trzy tryby:
\begin{itemize}
    \item \texttt{train} -- trening wybranego klasyfikatora,
    \item \texttt{eval} -- ewaluacja zapisanego modelu,
    \item \texttt{infer} -- predykcja z CSV lub kamery.
\end{itemize}

Przykładowe użycie:
\begin{lstlisting}[style=stylePython,caption={Przyklady uruchomien CLI}]
uv run python main.py train --classifier svc --model artifacts/svc_model.pkl
uv run python main.py eval --model artifacts/svc_model.pkl --dataset data/raw/sign_mnist_test.csv
uv run python main.py infer --source csv --row 42 --dataset data/raw/sign_mnist_test.csv
uv run python main.py infer --source camera --camera 0 --model artifacts/svc_model.pkl
\end{lstlisting}

\clearpage
\include{chapters/chapter007}
\clearpage
\chapter{Wnioski}
\label{cha:wnioski}

\section{Podsumowanie realizacji}
\label{sec:podsumowanie}

W ramach projektu zrealizowano kompletny pipeline klasyfikacji znaków dłoni:
od przygotowania danych tablicowych, przez uczenie i ewaluację modeli,
po inferencję z kamery w czasie rzeczywistym. Implementacja zachowuje
spójny format wejścia (28$\times$28, 784 cechy), co pozwala stosować
różne klasyfikatory bez zmian interfejsu danych.

\section{Ocena zastosowanych metod}
\label{sec:ocena_metod}

Przeprowadzone porównanie pokazuje, że:
\begin{itemize}
    \item modele klasyczne (\texttt{SVC}, \texttt{RandomForest}) zapewniają
          zadowalającą bazową skuteczność,
    \item model konwolucyjny (\texttt{CNN}) uzyskuje wyraźnie wyższą dokładność,
    \item jakość inferencji online zależy nie tylko od modelu, ale także od jakości
          detekcji i preprocesingu obrazu.
\end{itemize}

\section{Ograniczenia i ryzyka}
\label{sec:ograniczenia_ryzyka}

Najważniejsze ograniczenia obecnej wersji:
\begin{enumerate}
    \item klasyfikacja oparta na pojedynczych klatkach bez modelowania sekwencji,
    \item brak obsługi znaków dynamicznych wymagających ruchu,
    \item podatność na zmiany oświetlenia i tła,
    \item brak kalibracji pod indywidualne różnice użytkowników.
\end{enumerate}

\section{Kierunki dalszego rozwoju}
\label{sec:dalszy_rozwoj}

Rozszerzenia, które mogą poprawić jakość i użyteczność systemu:
\begin{itemize}
    \item dodanie modelu sekwencyjnego (np. LSTM/Transformer) dla rozpoznawania
          gestów dynamicznych,
    \item wygładzanie predykcji w oknie czasowym,
    \item augmentacja danych i rozszerzenie zbioru treningowego o rzeczywiste
          próbki z kamery,
    \item budowa modułu oceny wydajności czasu inferencji na różnych urządzeniach.
\end{itemize}

\section{Wniosek końcowy}
\label{sec:wniosek_koncowy}

Projekt spełnia cele laboratorium z zakresu sztucznej inteligencji:
łączy techniki uczenia maszynowego i przetwarzania obrazu w działającym
prototypie klasyfikatora znaków dłoni, który może stanowić bazę do
dalszych prac badawczych i inżynierskich.


% Wyłączenie działania ulem na czas bibliografii
\renewcommand{\emph}[1]{\textit{#1}}

\addcontentsline{toc}{section}{\textbf{Bibliografia}}
\bibliographystyle{plain}
\bibliography{bibliografia}

\renewcommand{\emph}[1]{\textit{#1}}

\clearpage
\addcontentsline{toc}{section}{\textbf{Spis rysunków}}
\listoffigures
\clearpage

\addcontentsline{toc}{section}{\textbf{Spis tabel}}
\listoftables
\clearpage

\addcontentsline{toc}{section}{\textbf{Spis listingów}}
\lstlistoflistings
\clearpage

\end{document}
